\documentclass[../main]{subfiles}
\begin{document}
\ifSubfilesClassLoaded{\setcounter{page}{50}}{}
\section{Communism and the Family}
\label{sec:09_communism_family}

“According to the materialist conception, the determining factor in history is, in the last instance, the production and reproduction of immediate life. This, again, is of a twofold character. On the one side, the production of the means of subsistence: of articles of food, clothing and shelter, and the tools necessary for that production; on the other side, the production of human beings themselves, the propagation of the species. The social organisation under which the people of a particular historical epoch and a particular country live is determined by both kinds of production: by the stage of development of labour on the one hand, and of the family on the other,” wrote Frederick Engels in The Origin of the Family, Private Property and the State.

In the previous chapters, we discussed how, under communism, the production of the necessities of life is automated and subordinated to the production of human life in all its fullness. Society consciously creates its own development and its own forms of collectivity, expanding the content of free activity. The family, as an economic unit based on the private character of appropriation of the products of labour, loses its significance and disappears during the period of transition. “Private housekeeping is transformed into a social industry,” wrote Engels. We saw above how much social time this saves, freeing people from the burdens of everyday life. However, socialisation concerns more than just consumption.

The main revolution in family relations lies in the social upbringing of children. This revolution has long been brewing, and movement in this direction already exists under capitalism. Housekeeping has begun to be “socialised”, and everyday life has long since become, in many respects, a branch of the social economy. We cook mostly from semi-prepared foods, buy our clothes in shops and wash them at the laundry, even though not long ago all this was women’s work done at home. The modern man, too, no longer needs to carry water, chop firewood or build the house his family lives in. All of this has long been a high-tech social undertaking.

But the upbringing of children still refuses to become a sphere of production, even though public institutions for this purpose appeared long ago and have been actively developing for at least two centuries. And the better-off have always sent their children precisely there and hired private tutors, while they themselves merely enjoyed the comforts of family life. So it was in the Soviet Union, too: one of the first achievements that socialist working people spoke of with such pride was the organisation of factory nurseries and kindergartens. In other words, the need to entrust the upbringing of children to those who know how to do it, and to win free time for what one considers truly necessary and important (for some, incidentally, that may be the upbringing of children—why not?), is something rich and poor alike share, with only a few differences of a petty-bourgeois, sentimental nature. Capitalism has long been destroying the traditional family, and this process has no moral roots, only economic ones. At present, the promotion of family values is driven, on the one hand, by the achievements of socialist revolutions, which banned the exploitation of child labour, and on the other, by capital’s need for a market in children’s commodities and education. But there is no real historical or social necessity to preserve the family in its current form as an economic unit.

Thus, as society transitions to communist relations, the sphere of the upbringing of children will become a sphere of social production. Upbringing becomes genuine social creativity—free activity by the standards of truth, goodness and beauty. This increases many times over the responsibility of educators both to children and to society as a whole. As A. S. Makarenko once wrote: “We must have no childhood catastrophes, no failures, no percentage of rejects, even if expressed in hundredths of a unit!” In a number of works, he compared upbringing with other branches and showed how much more dangerous failure is here than anywhere else. Speaking of the family, Anton Semyonovich observed that the task of parents is to raise worthy members of society; they are accountable to society for the “quality of the material supplied”, and in return they receive the warmth of family relationships. But by no means all parents know how to raise children. Moreover, even now, by no means all children have parents or simply a complete family. Herein lies the need for a conscious, high-quality socialisation of upbringing and for the creation of other forms of “family”.

The communist “family” is a collective form in which everyone takes part in producing social life on a correspondingly social scale. It follows that the size of such a “family” may vary, and it certainly cannot be confined within its present-day bounds. Its size depends on many concrete-historical conditions of social production, so various forms will be possible, both within a single epoch in the development of communist society and over the course of its history. Since such a “family”, as the primary collective of the human being in the making, will be founded on the social character of activity and directly determined by it, the principle on which such collectives are formed will likewise be directly social—qualitatively different from what it is now. What is at stake is the abolition of modern marriage as a form, especially since sexual relations and sexual love, on the one hand, and marital-economic relations, on the other, are far from being one and the same thing. That is why one of A. S. Makarenko’s most important and enduring practical discoveries was that the new “family” must draw every person directly into social production as an active member of the social labour collective, rather than as a consumer of social goods or as an “object” of upbringing.

Tracing the rapid, at times barely perceptible changes in this sphere during the very first years of the Great October Socialist Revolution, Alexandra Kollontai, in her work “Make Way for Winged Eros!”, sought to discern and predict the type of personal relations, or more precisely relations of sexual love, that would take hold once social property had completely triumphed over private property, as new social institutions developed and as elements of an infrastructure capable of relieving the burden of motherhood were introduced, increasingly turning it into a branch of social production.

Wherever the new everyday life was established (and not only in the USSR), it truly emancipated women on a scale unimaginable to the rest of the world, giving them time not only for love itself and for “courtship”, but also for rethinking love. These new, already partly communist relations are also depicted in the fiction of the industrialisation era, both in the prose of Alexandra Kollontai herself (“Vasilisa Malygina”, “A Great Love”) and, for example, in Ilya Ehrenburg’s novel “The Second Day”.

Establishing the commune without the registration of marriage (in the 1920s and 1930s, unregistered marriage became the norm, while registration, on the contrary, came to be seen as a sign of mistrust between partners) where the patriarchal family had previously reigned was, of course, getting far ahead of events. But it was a revolution in how lovers, those in love, married people and divorcees perceived themselves. Not marriage as a life for two and for the sake of two, but family life directly shaped by the goals of a rapidly growing and developing society and of one’s immediate (labour) collective—that is the real agenda of love, one that expresses the tendency towards communism.

In capitalist society, the family, as a zone of love and of the responsibility bound up with it, is a “high-risk zone”. Within the walls of their home, spouses are doomed to live through everything that is programmed by the relations of private property, by questions of means, earnings, spending and the rest of life’s prose. But communism is precisely the elimination of such problems.

The dynamics of the withering away of both the old family and the state as such are complex. However, the half-measures of socialism are powerless here: students and young specialists, leaving the commune-hostels, still ended up in flats, in families, within the old frameworks, only now with new social security. The motivation to outgrow the old cells of the “honeycomb” was lost in a society where commodity-money relations were once again winning back their dominant position. The leap towards the commune was followed by a natural, and seemingly also temporary, return to patriarchy. Once again, it was the care of the parents alone, rather than of the whole collective or society, that came to the fore in upbringing. And yet what was emerging was not only “sons of the regiment” (the war had its own say in this), but also children of the communes.

It was in these “children of the communes” that one of the most important aspects of the becoming of communism was embodied. In the transition to communism, children must cease to be divided into one’s own and other people’s. There are no other people’s children (that is, children who are someone else’s responsibility to care for, not mine) in a vast collective where the problems of the “meter” measuring one’s personal labour contribution have been removed.

Finding themselves from the very start in a collective with many men and women, children, while still distinguishing their father and mother, will not see the “edge and end of the world” at the point where mum and dad end and the “neighbours” begin. Through the socialisation that comes naturally in a commune, through the everyday experience of the collective’s care, children will stop “finishing their games in the sandbox of private property”. Moving away from the old form of marriage will make it possible to avoid a number of vestiges reproduced in the psyche and in the consciousness of fathers and mothers, and not only of children. In a large and close-knit collective, a change of partners ceases to be something traumatic if there is no longer love between the members of a couple. In general, in the absence of social economic barriers, everything strictly personal (including the intimate) can only then become truly personal, rather than opposed to the social, once freed from economic, legal and moral shackles.

The new “family” must itself become a directly collective form of communist activity aimed at creating the new man, one that directly links all its members, including children, with society as active subjects.
\end{document}
